% R32 section: follow-up (2026-10-15) -- the open points, answered.  Input by R32_floquet_spectra.tex.
% Every statement here is pinned by a test in code/qca_fragmentation/tests/test_spectral.py.
\section{Exact structures behind the classes}
\label{sec:answers}
This section collects the exact structures found in the follow-up (2026-10-15), which explain the classes: the flip map,
the hidden time reversal of the one-sided rules, the universal scar, the kink dispersion of 156/198 and the domain-wall
charge of X-57. It also reports the block-first runs beyond the dense-sector limit (step~9), the eigenvector analyses at
small $N$ (step~10), and the checks of the assumptions those runs make. Every statement is pinned by a test in
\fn{code/qca\_fragmentation/tests/test\_spectral.py}.

\subsection{Flip partners: the X flip turns $H$ into $XHX$}
The physical spin flip is $X^{\otimes N}$. It maps a rule $W$ onto its partner $W'$ (the firing pattern
$f(m,n)\to f(1-m,1-n)$), but the gate changes with it:
\begin{equation}
  X^{\otimes N}\,U_W(H)\,X^{\otimes N}=U_{W'}(XHX),\qquad XHX=\tfrac{1}{\sqrt2}(X-Z)=-ZHZ .
\end{equation}
Because $Z^{\otimes N}$ commutes with the controls, $Z^{\otimes N}U_{W'}(G)Z^{\otimes N}=U_{W'}(ZGZ)$. So
$U_W(H)\cong U_{W'}(-H)$ up to the basis change $Z^{\otimes N}$: the partner rule with a $-1$ on every firing. This is a
conditional phase, not a global one. Verified as operators at $N=8$, on the ring, for the pairs 54/147, 108/201, 60/195,
102/153 and 57/99. The consequences are:
\begin{itemize}
\item On the ring, flip partners have \emph{identical} supports, hence identical Krylov sectors and blocks, but
  \emph{different} spectra. Examples: 108 and 201 ($\langle r\rangle=0.399$ and $0.417$ from identical block structures),
  60/102 and 153/195 (0.511 and 0.531), and 54 and 147 (0.483 and 0.518).
\item The obc0 boundary (ghost sites fixed to $0$) breaks the flip already for $X$. With $X$ on the ring, 108 and 201 have
  4156 cycles each; on the open chain they have 4868 and 3492.
\item 57/99 and 156/198 are also \emph{reflection} partners, and the reflection is an exact symmetry with either gate.
  Their spectra are identical in every dataset.
\end{itemize}

\subsection{The hidden symmetry of the one-sided rules}
On the ring, the one-sided rules have $\beta=1$ statistics without mirror symmetry (Sec.~\ref{sec:pbc}). The
time-reversal symmetry behind this is now identified:
\begin{equation}
  S=B_e\,D\,R_\mathrm{bond}\,B_e^{\mathsf T},\qquad A=K\,L_A\,S,
\end{equation}
where:
\begin{itemize}
\item $R_\mathrm{bond}$ is the bond-centred reflection $i\to1-i$, which exchanges the two sublattices;
\item $B_e$ rotates the even sites into the eigenbasis $\{|h_+\rangle,|h_-\rangle\}$ of the Hadamard;
\item $D=X^{\otimes N}$ for 153/195, and a diagonal sign operator for 60/102.
\end{itemize}
Verified as explicit operators at $N=8$ and $10$, for all four rules, to $10^{-13}$:
\begin{itemize}
\item $S$ is real, $[S,U]=0$, and $ST_2S^{-1}=T_2^{-1}$ ($S$ flips the momentum);
\item $A$ preserves the momentum and maps $U\to U^{-1}$ and $\Omega_T\to\Omega_T^{-1}$;
\item $A^2=+1$: COE, not Kramers pairs.
\end{itemize}
In the mixed frame (Hadamard basis on the targets of one layer, computational basis on its controls) that layer is
diagonal, and $S$ becomes a dressed reflection. This is why a search in the computational basis could not find it. The
single-site rotation needed is $V=(|h_+\rangle,|h_-\rangle)$, a rotation by $\pi/8$, not $H$. For 57/99 the same search
returns nothing, consistent with CUE.

\paragraph{Correction to Sec.~\ref{sec:pbc}.} In the real momentum sectors ($j=0$, $N/4$), $S$ preserves $j$ and commutes
with $U$. So the real sectors of the one-sided rules carry an unresolved $\mathbb Z_2$, exactly as the reflection does for
the mirror rules. Their real-sector values (0.414 and 0.428 at $N=14$) mix two symmetry sectors and are not
diagnostic. The conclusions used only the complex sectors, so they are unchanged.

\subsection{Scars: one universal eigenstate, edge-projected}
\textbf{$|h_+\rangle^{\otimes N}$ is a $+1$ eigenvector of every controlled-Hadamard gate, for every rule and boundary},
because $H|h_+\rangle=|h_+\rangle$ whatever the controls read. All gate singlets of the chaotic rules follow from it
(found from the gate algebra at $N\le12$ and identified site by site):
\begin{itemize}
\item \emph{$\theta=0$:} $|h_+\rangle^{\otimes N}$ itself.
\item \emph{Rules with a frozen vacuum} (54 on both boundaries, 60/102 on the ring): its projection off the vacuum,
  $|\psi\rangle\propto|h_+\rangle^{\otimes N}-\cos^N(\pi/8)\,|0\ldots0\rangle$. This is the ``entangled singlet''. It has
  Schmidt rank 2, its entanglement follows in closed form and matches the data to $10^{-6}$ for $N=6\ldots14$, and it
  falls to zero as $\cos^N(\pi/8)\to0$. It is the same state in all three rules (overlap $1$).
\item \emph{Boundary singlets} ($\theta=\pi$, or $0$ with two of them): $|h_-\rangle$ on an end site whose gate fires
  for both values of the inner neighbour. With the ghost fixed to $0$, that site picks up $-1$ per period
  deterministically. This predicts the counts exactly: 57 (right end) 2, 99 (left end) 2, 147 (both ends) 4, 54 (neither
  end, vacuum frozen) 1. On the ring there are no ends, hence exactly one singlet.
\end{itemize}
These are exact scars, but of the simplest kind: a product eigenstate of the gate (or its projection) embedded in a
chaotic spectrum, in the spirit of the projector embedding~\cite{shiraishi2017}.

\paragraph{The outliers of 57.} These are \emph{not} singlets. They are an \emph{edge mode on the
$|h_+\rangle^{\otimes N}$ vacuum}: away from the left end the state is $|h_+\rangle$ on every site (overlap rising from
0.86 to 0.93 into the bulk), and the disturbance is localised on the first four sites. The two quasienergies are one
mode: $\pm0.3122\pi$ alone, and $\pm0.6877\pi=\pm0.3122\pi+\pi$ together with the right-end $|h_-\rangle$ singlet. The mode
lives at the left end because there 57's gate depends on the inner neighbour (it fires on $00$, not on $01$); the right
end is inert. On the ring the strongly non-thermal states disappear (minimum $\eta$ 0.41 to 0.55, against 0.13 to 0.27
on the open chain).

\subsection{156/198: one kink per segment, with an exact dispersion}
For 156 (fires only on left $=1$, right $=0$), a $01$ domain wall is frozen and a $10$ wall is mobile. Between two frozen
walls a segment of $L$ free sites holds exactly one kink, in $L+1$ positions. No gate couples two segments, so
$U=\bigotimes_s U_s$. The single-kink spectrum is
\begin{equation}
  \cos\theta_m=-\sin^2\!\frac{\pi m}{L+1},\quad m=1\ldots\lfloor(L+1)/2\rfloor\ (\pm\text{ pairs; }\theta=\pi\text{ at }m=\tfrac{L+1}{2}),
  \qquad\text{plus }\theta=0,
\end{equation}
i.e.\ $\cos\theta(q)=(\cos q-1)/2$ with $q=2\pi m/(L+1)$. It does not depend on the brickwork offset, and it reproduces
the exact spectra of the largest Krylov sectors at $N=12$ as sums over their segments (including the end segments).
Blocks have at most about 125 states, so their levels are sums of a handful of values from this small set. That gives
many exact and near coincidences, hence $\langle r\rangle\approx0.42$: neither independent levels nor level repulsion.

\subsection{X-57: the ergodic orbits need the boundary}
On the ring, X-57 (and X-99) \emph{conserves the number of domain walls}. A search over sublattice-resolved 3- and
4-site densities finds exactly this one local charge (plus its 4-site extension) on the ring, and none on the open
chain. There the ghost sites create and annihilate walls. Ring orbits are confined to fixed wall number, so their share
of $2^N$ decays (0.007 at $N=14$). On the open chain the boundary mixes the wall sectors and the orbits fill a finite
fraction. For contrast, 54 and 150 keep local charges on both boundaries.

\IfFileExists{R32_answers_step9.tex}{% R32: beyond the dense-sector limit (step 9), eigenvectors at small N (step 10), and the checks of their assumptions.
% Data: results_spectra/blockfirst/{ring,tower,verify}_*; code: spectral/blockfirst.py; tests: test_blockfirst_matches_pipeline.
\subsection{Beyond the dense-sector limit (block-first)}
The production pipeline diagonalises whole Krylov sectors densely, so its memory is set by the largest \emph{sector}
(Sec.~\ref{sec:memory}). For the cases that matter for the open questions, an exactly known symmetry was used to split
the space \emph{before} the dense step, by a sparse construction (\fn{spectral/blockfirst.py}), and only eigenvalues were
computed.
\begin{itemize}
\item \emph{Ring, $N=16$} (57, 54, 147, 153, 60): $T_2$-momentum orbit bases first ($\approx8200$ states per sector),
  then the half step $\Omega_T$ inside each sector. The eigenphases come from a Cayley transform, a Hermitian
  eigenvalue problem.
\item \emph{Open chain, one-sided rules:} the largest tower block, built directly in its product basis. For 153 it is
  $|h_-\rangle$ on site 0 with the other sites free ($N=15$); for 60 it is site $0=|1\rangle$ and site $N-1=|h_-\rangle$
  ($N=16$). Each has $16\,384$ states.
\end{itemize}
Validation against the production data: the ring code reproduces the Julia $\varphi$ of every momentum sector to
$10^{-14}$ ($N=12$) and the per-block $\langle r\rangle$ to $10^{-11}$ ($N=12$, 14). The tower blocks reproduce the
Julia largest blocks level by level to $10^{-13}$ ($N=10\ldots13$).

\emph{One pitfall, caught by that comparison:} computing $|\theta|$ from $\arccos$ of the eigenvalues of
$(U+U^{\mathsf T})/2$ smears exactly degenerate eigenvalues $\pm1$ by $\sqrt{\epsilon}\approx10^{-8}$ into spurious
near-duplicates. The real tower blocks therefore use the eigenvalues of $U$ directly.

\emph{Results} (Secs.~\ref{sec:pbc} and \ref{sec:results}):
\begin{itemize}
\item On the ring at $N=16$: 57 is $0.5985\pm0.0011$ (CUE), 147 is $0.528$ and 153 is $0.527$ (COE-shaped), and 60
  ($0.523$) and 54 ($0.510$) are still rising, now COE-shaped.
\item On the open chain the one-sided statistics stop drifting: $\langle r\rangle\approx0.424$ for 153 at $N=15$ and for
  60 at $N=16$.
\end{itemize}

\paragraph{Checking the assumptions (user request).} Block-first assumes that nothing splits the pieces further. This
was tested at the production sizes with the generic-$K$ trick:
\begin{itemize}
\item \emph{Ring:} restricted to a momentum sector, $K$ is a generic combination of $T_2$-invariant elements of the gate
  algebra ($S_A$, $S_B$, $\{S_A,S_B\}$, $U+U^\dagger$). One component per Krylov sector means that no hidden symmetry
  splits the sector: any such symmetry would also split this subalgebra.
\item \emph{Tower:} $K=\sum_ia_ig_i$ restricted to the block.
\end{itemize}
Results:
\begin{itemize}
\item In every tested complex momentum sector at $N=16$ (57 $j{=}1$; 54, 153 $j{=}1,2$; 147 $j{=}1$; 60 $j{=}1,3$) the
  test gives exactly one component, equal to the single Krylov component. $K$ is non-degenerate (gaps
  $4\times10^{-6}$ to $2\times10^{-4}$), and the sector is invariant under $\Omega_T$ to $2\times10^{-15}$.
\item The tower block of 153 at $N=15$ is irreducible: one component of $16\,384$, gap $2.3\times10^{-6}$, and the gates
  leave the block invariant to $1.4\times10^{-14}$.
\item The tower block of 60 at $N=16$ is irreducible: one component of $16\,384$, gap $2.3\times10^{-7}$, and gate
  invariance $2\times10^{-14}$.
\end{itemize}
In the \emph{real} momentum sectors the same test splits the space into two halves. This is how the residual
$\mathbb Z_2$ of the one-sided rules was found, which led to their time reversal (above). The conclusions use only the
complex sectors.

\paragraph{Eigenvectors at small $N$ (step 10).} The scar analyses use eigenvectors from the exact engine at
$N\le12$:
\begin{itemize}
\item the gate singlets, from the one-dimensional gate-algebra blocks;
\item the outliers of 57 and the ETH matrix elements, from the dense eigenvectors of $U$ in the largest block.
\end{itemize}
The Julia pipeline was not changed, so no full recompute was needed.
}{}
